\documentclass[tikz,border=6pt]{standalone}
\usepackage{ctex}
\usepackage{amsmath}
\usepackage{pgfplots}
\pgfplotsset{compat=1.18}
\usetikzlibrary{calc,angles,quotes,arrows.meta,positioning,decorations.markings}
\begin{document}
\begin{tikzpicture}[scale=4.2]
  % 坐标轴
  \draw[->,gray!70,thick] (-1.35,0) -- (1.45,0) node[right,black]{$x$};
  \draw[->,gray!70,thick] (0,-1.35) -- (0,1.45) node[above,black]{$y$};

  % 单位圆
  \draw[blue!60!black,very thick] (0,0) circle (1);
  \node[blue!55!black,font=\small] at (-0.5,-0.6) {单位圆 $r=1$};

  % 一周特殊角的终边与圆上点
  \foreach \deg in {0,30,45,60,90,120,135,150,180,210,225,240,270,300,315,330}
  {
    % 终边
    \draw[gray!55,thin] (0,0) -- (\deg:1);
    % 圆上点
    \fill[red!80!black] (\deg:1) circle (0.012);
  }

  % 突出展示四个一象限特殊角的弧度与坐标（避免整圈拥挤）
  \foreach \deg/\rad in {0/$0$,30/$\frac{\pi}{6}$,45/$\frac{\pi}{4}$,60/$\frac{\pi}{3}$,90/$\frac{\pi}{2}$,
    120/$\frac{2\pi}{3}$,135/$\frac{3\pi}{4}$,150/$\frac{5\pi}{6}$,180/$\pi$,
    210/$\frac{7\pi}{6}$,225/$\frac{5\pi}{4}$,240/$\frac{4\pi}{3}$,270/$\frac{3\pi}{2}$,
    300/$\frac{5\pi}{3}$,315/$\frac{7\pi}{4}$,330/$\frac{11\pi}{6}$}
  {
    \node[font=\scriptsize,red!65!black] at (\deg:1.13) {\rad};
    \node[font=\tiny,gray!75!black] at (\deg:1.27) {$\deg^\circ$};
  }

  % 重点标注：theta=60° 的点 P=(cos,sin)
  \coordinate (P) at (60:1);
  \fill[red!80!black] (P) circle (0.02);
  % 投影线
  \draw[dashed,green!45!black] (P) -- (0.5,0);
  \draw[dashed,orange!70!black] (P) -- (0,0.866);
  \node[green!40!black,font=\scriptsize] at (0.25,-0.07) {$\cos\theta$};
  \node[orange!65!black,font=\scriptsize,rotate=90] at (-0.07,0.45) {$\sin\theta$};

  % 角弧
  \draw[->,purple!70,thick] (0.32,0) arc (0:60:0.32);
  \node[purple!70,font=\small] at (28:0.43) {$\theta$};

  % P 坐标说明框（半透明，置于第一象限空白处偏外）
  \node[draw=red!50!black,fill=white,fill opacity=0.85,text opacity=1,
        rounded corners,font=\scriptsize,align=left] at (0.62,1.18)
    {终边交点\\ $P=(\cos\theta,\sin\theta)$};
\end{tikzpicture}
\end{document}
